\section{Introduction}
\label{sec:introduction}

The landscape of software development is undergoing a profound transformation. In 2023, GitHub reported that in files where Copilot is enabled, 46\% of code is now written with AI assistance~\citep{github2023copilot}. Tools like GitHub Copilot have surpassed one million paid subscribers, while autonomous coding agents like Devin and SWE-agent have demonstrated the ability to independently resolve real-world software issues. The emergence of ``vibe coding''---where developers describe intent and iterate with AI rather than writing code directly---signals a fundamental shift in how software is created.

Yet this rapid evolution has created significant challenges for researchers and practitioners. The field lacks a unified framework for understanding the capabilities and limitations of different systems. Questions abound: How do code completion tools differ from conversational assistants? What makes an ``agent'' autonomous? Which benchmarks meaningfully measure real-world utility? Where should research efforts focus?

\subsection{Contributions}

This survey addresses these challenges by providing:

\begin{itemize}
    \item \textbf{A comprehensive taxonomy} categorizing coding agents along three dimensions: autonomy level, interaction scope, and integration mode.
    \item \textbf{Systematic coverage} of major systems from the LLM era (2022--present), including code completion tools, conversational assistants, autonomous agents, and integrated development environments.
    \item \textbf{Critical analysis} of evaluation benchmarks, examining what they measure and what they miss.
    \item \textbf{Identification of open challenges} and a research agenda for the software engineering community.
\end{itemize}

\subsection{Scope}

We focus on tools that assist with code writing and modification, concentrating on systems from the LLM era (2022--present) that have publicly documented capabilities and evaluation data. We exclude general program synthesis, formal verification tools, and non-LLM static analysis systems.

\subsection{Organization}

The remainder of this paper is organized as follows. Section~\ref{sec:background} provides background on large language models for code and the agent paradigm. Section~\ref{sec:taxonomy} presents our taxonomy of coding agents. Sections~\ref{sec:completion}--\ref{sec:ide} survey four categories of tools in detail. Section~\ref{sec:benchmarks} examines evaluation benchmarks and metrics. Section~\ref{sec:challenges} discusses open challenges and future directions. Section~\ref{sec:conclusion} concludes.
